% ============================================================================
% Fichier  : partie4/P04_C10_identifier.tex
% Titre    : Identifier les Vulnérabilités sans les Exploiter
% Auteur   : MINKA MI NGUIDJOÏ Thierry Emmanuel
% Version  : 1.0 -- novembre 2024
% Corrigés : corriges/Corr_C10.tex
% ============================================================================
% Sources : NIST SP 800-115 (techniques d'évaluation) ; NIST SP 800-40 r4
% (gestion des correctifs) ; CIS Benchmarks ; loi n° 2010/012 (art. 65 al.
% 2 et 69, rappel).
% Faits nouveaux (bible) : SOLITRANS, inventaire CMDB 1 380 actifs, 1 512
% découverts (132 inconnus) ; balayage authentifié réussi sur 78 % des
% serveurs ; 3 812 constats bruts -> 2 947 après dédoublonnage ; échantillon
% de 60 constats vérifiés : 9 faux positifs ; incident du 18/05/2025 sur un
% terminal de pesage à Bonabéri. NYANGA : liste de cibles tirée de la CMDB,
% passerelle absente.
% ============================================================================

\chapter{Identifier les Vulnérabilités sans les Exploiter}
\label{chap:identifier}
\chaptermark{Identifier sans exploiter}

\epigraph{%
    « Un outil de balayage ne trouve que ce qu'on lui montre,\\
    et il ne montre que ce qu'il a su lire. »%
}{Principe du cours}

\niveauChapitre{Intermédiaire}{%
    Maîtrise du Chapitre~\ref{chap:vulnerabilite}
}

\begin{boxobjectifs}
\begin{itemize}
    \item Construire et tenir l'inventaire des actifs, condition de toute
          identification.
    \item Choisir entre balayage non authentifié, authentifié, agent et
          analyse de configuration.
    \item Reconnaître les faux positifs et les faux négatifs, et en
          connaître les causes.
    \item Encadrer un balayage~: autorisation, périmètre, fenêtre, effets
          de bord.
    \item Mesurer la couverture d'un programme d'identification.
\end{itemize}
\end{boxobjectifs}

\bigskip

% ============================================================================
\section*{Avant de commencer : la liste des cibles}
% ============================================================================

Bafoussam, avril 2025. Mme~NDOUMBE cherche pourquoi l'outil de balayage de
NYANGA CRÉDIT n'a jamais signalé la faille de la passerelle. L'outil était
acheté, configuré et lancé chaque mois. Elle ouvre sa configuration~: la
liste des adresses à analyser est importée de la base de gestion des
configurations, tenue par le DSI. La passerelle n'y figure pas. Elle a été
installée en 2023 par un prestataire, dans une zone réseau à part, et
personne ne l'a jamais déclarée. L'outil a fait exactement ce qu'on lui
demandait.

\begin{boxalert}
\textbf{Violation typique.} Mesurer un programme de vulnérabilités au
nombre de failles trouvées, sans mesurer la part du parc réellement
analysée.

\smallskip
\textbf{Règle~:} l'identification commence par l'inventaire et se juge à sa
couverture. Un actif absent de l'inventaire, ou analysé sans succès, est un
\NEK{}~: on ne sait rien de ses failles.
\end{boxalert}

\separateur

% ============================================================================
\section{L'Inventaire d'Abord}
\label{sec:c10-inventaire}
% ============================================================================

On ne corrige que ce que l'on connaît. L'inventaire des actifs (matériels,
systèmes, logiciels et services) est le premier contrôle de tout programme,
et le plus souvent défaillant. Il se construit en croisant plusieurs
sources, parce qu'aucune n'est complète.

\begin{tableaucours}{L{3.6cm}L{4.4cm}L{4.2cm}}{%
    \tableauHeader{Source} &
    \tableauHeader{Ce qu'elle voit} &
    \tableauHeader{Ce qu'elle manque}
}
Base de gestion des configurations & Ce qui a été déclaré & Ce qui a été
installé sans déclaration \\
Annuaire (Active Directory) & Machines Windows du domaine & Équipements
réseau, Linux, appareils hors domaine \\
Serveur DHCP & Appareils qui demandent une adresse & Adresses fixes \\
Découverte réseau (balayage des plages) & Tout ce qui répond & Ce qui est
éteint, filtré ou hors des plages \\
Console des fournisseurs en nuage & Ressources en nuage & Ce qui est
ailleurs \\
\end{tableaucours}
\vspace{-0.8em}
\begin{center}{\small\itshape Tableau~10.1. Les sources de l'inventaire}\end{center}

\begin{boxscenario}{Opération SOLITRANS. Les 132 inconnus}
\textbf{Constatations.} La base de gestion des configurations déclare
1\,380 actifs. La découverte réseau des plages du siège, de Bonabéri et de
Kribi, croisée avec l'annuaire et le serveur DHCP, en trouve 1\,512. Les
132 actifs non déclarés comprennent 41 caméras et enregistreurs vidéo, 28
terminaux de pesage et de lecture de conteneurs, 19 postes de prestataires,
17 imprimantes, 14 commutateurs et 13 appareils non identifiés.

\textbf{Analyse} (confiance élevée). Un actif sur onze échappait à toute
analyse. Les catégories manquantes sont celles qu'installent les
prestataires ou les services métier, comme la passerelle de NYANGA.
\end{boxscenario}

% ============================================================================
\section{Quatre Façons d'Identifier}
\label{sec:c10-methodes}
% ============================================================================

\begin{tableaucours}{L{2.8cm}L{3.8cm}L{2.8cm}L{2.8cm}}{%
    \tableauHeader{Méthode} &
    \tableauHeader{Principe} &
    \tableauHeader{Force} &
    \tableauHeader{Faiblesse}
}
Balayage non authentifié & Interroger les services ouverts depuis le
réseau, lire bannières et réponses & Voit ce que voit un attaquant &
Déduit les versions, faux positifs nombreux \\
Balayage authentifié & Se connecter avec un compte en lecture et lire les
paquets installés & Précis sur les versions et correctifs & Exige un
compte~; échecs de connexion silencieux \\
Agent & Programme installé sur la machine, qui remonte son état & Couvre les
portables hors réseau & À déployer et maintenir \\
Analyse de configuration & Comparer la configuration à une référence (CIS
Benchmarks) & Trouve les réglages dangereux, sans CVE & Ne voit pas les
failles logicielles \\
\end{tableaucours}
\vspace{-0.8em}
\begin{center}{\small\itshape Tableau~10.2. Quatre méthodes d'identification}\end{center}

Aucune de ces méthodes n'exploite la faille. Elles constatent qu'un logiciel
vulnérable est présent, ou qu'un réglage est dangereux. Démontrer qu'un
attaquant peut s'en servir relève du test d'intrusion, conduit dans un
autre cadre (UE MSN-2S3-SEC03). Le guide NIST SP~800-115~\cite{nist800115}
décrit les deux familles et leurs précautions.

\begin{boxinfo}
\textbf{Les composants logiciels.} Une application peut embarquer une
bibliothèque vulnérable sans que le système d'exploitation le sache. Les
outils d'analyse de composants lisent les fichiers de l'application, ou une
liste de ses composants fournie par l'éditeur, pour retrouver ces
bibliothèques. Un balayage classique ne les voit pas.
\end{boxinfo}

% ============================================================================
\section{Faux Positifs, Faux Négatifs}
\label{sec:c10-faux}
% ============================================================================

\begin{boxdef}
\textbf{\termecle{Faux positif}}~: vulnérabilité signalée alors qu'elle est
absente. \textbf{\termecle{Faux négatif}}~: vulnérabilité présente et non
signalée.
\end{boxdef}

Les faux positifs coûtent du temps et de la crédibilité~: une équipe système
qui corrige trois failles inexistantes cesse de croire le rapport. Les faux
négatifs coûtent plus cher, parce qu'ils ne se voient pas.

\begin{tableaucours}{L{6.0cm}L{6.2cm}}{%
    \tableauHeader{Cause de faux positif} &
    \tableauHeader{Cause de faux négatif}
}
Correctif rétroporté~: la distribution corrige sans changer le numéro de
version & Actif absent de la liste des cibles \\
Bannière de service modifiée ou générique & Authentification échouée,
analyse réduite au non authentifié \\
Composant présent mais non utilisé & Service filtré par un pare-feu entre
l'outil et la cible \\
Règle de détection trop large & Règle de détection absente pour la faille
ou le produit \\
\end{tableaucours}
\vspace{-0.8em}
\begin{center}{\small\itshape Tableau~10.3. Les causes d'erreur d'un balayage}\end{center}

\begin{boxinfo}
\textbf{Vérifier un constat avant de le transmettre}
\begin{enumerate}
    \item Relire la preuve donnée par l'outil~: version lue, réponse du
          service, paquet installé.
    \item Confirmer la version par une autre voie (connexion à la machine,
          console d'administration).
    \item Vérifier les bulletins de la distribution pour un éventuel
          correctif rétroporté.
    \item Confirmer que le composant vulnérable est actif.
    \item Coter le constat~: confirmé, faux positif, ou non vérifiable.
\end{enumerate}
\end{boxinfo}

\begin{boxscenario}{Opération SOLITRANS. Soixante constats vérifiés}
\textbf{Constatations.} Après dédoublonnage (une même faille signalée par
plusieurs règles, ou sur plusieurs ports d'un même actif), les 3\,812
constats bruts deviennent 2\,947. Mme~EBODE fait vérifier un échantillon de
60 constats tirés au hasard~: 46 sont confirmés, 9 sont des faux positifs
(dont 7 correctifs rétroportés sur des serveurs Linux), 5 ne sont pas
vérifiables faute d'accès à la machine. Par ailleurs, l'authentification a
réussi sur 78~\% des serveurs~; sur les autres, l'outil s'est contenté d'un
balayage non authentifié sans le signaler dans le résumé.

\textbf{Analyse} (confiance moyenne, l'échantillon étant petit). Environ
15~\% des constats seraient des faux positifs. Et 22~\% des serveurs ont
été analysés de façon superficielle~: le nombre de failles réelles sur ces
serveurs est \NEK{}, sans doute supérieur à ce que dit le rapport.
\end{boxscenario}

% ============================================================================
\section{Encadrer un Balayage}
\label{sec:c10-encadrer}
% ============================================================================

Un balayage interroge des machines, parfois des milliers de fois. Il peut
ralentir un service, saturer un équipement fragile ou déclencher des
alertes. Et il touche au droit~: analyser un système sans autorisation peut
constituer un accès non autorisé (loi n°~2010/012, art.~65 al.~2, voir
Chapitre~\ref{chap:normes-droit}).

\begin{tableaucours}{L{3.4cm}L{8.8cm}}{%
    \tableauHeader{Élément} &
    \tableauHeader{Contenu}
}
Autorisation écrite & Signée par le propriétaire des actifs, et par
l'hébergeur ou le prestataire quand l'actif n'appartient pas à
l'organisation \\
Périmètre & Plages d'adresses, exclusions motivées, actifs fragiles \\
Fenêtre & Jours et heures, compatibles avec l'activité \\
Profil & Intensité, tests exclus (déni de service), comptes utilisés \\
Contacts & Personne à joindre en cas d'effet de bord, de part et d'autre \\
Information & Équipe de supervision prévenue, pour ne pas confondre le
balayage avec une attaque \\
\end{tableaucours}
\vspace{-0.8em}
\begin{center}{\small\itshape Tableau~10.4. Le cadre d'un balayage}\end{center}

\begin{boxscenario}{Opération SOLITRANS. Le terminal de pesage}
\textbf{Constatations.} Le 18~mai 2025 à 10~h~20, pendant le balayage du
site de Bonabéri, un terminal de pesage des conteneurs cesse de répondre. La
file des camions s'allonge pendant quarante minutes, jusqu'au redémarrage
manuel. Le terminal, non déclaré dans la base, n'était pas exclu du
périmètre.

\textbf{Analyse} (confiance élevée). Les équipements industriels et
portuaires supportent mal les balayages agressifs. Ils doivent être
inventoriés, exclus du balayage actif ou analysés avec un profil léger, en
dehors des heures d'exploitation.
\end{boxscenario}

% ============================================================================
\section{Mesurer la Couverture}
\label{sec:c10-couverture}
% ============================================================================

\begin{boxlecture}{La couverture de SOLITRANS en mai 2025}
\begin{tableaucours}{L{7.4cm}cc}{%
    \tableauHeader{Question} &
    \tableauHeader{Valeur} &
    \tableauHeader{Cote}
}
Tous les actifs sont inventoriés & 1\,380 sur 1\,512 & \NON{} \\
Tous les actifs inventoriés ont été analysés & 1\,342 sur 1\,380 & \NON{} \\
Les serveurs ont été analysés avec authentification & 78~\% & \NON{} \\
Les constats sont vérifiés & Échantillon de 60 & \NEK{} \\
Les actifs fragiles sont protégés des effets de bord & Incident du 18~mai &
\NON{} \\
\end{tableaucours}

\medskip
\textbf{Lecture.} Quatre \NON{} et un \NEK{}~: le premier balayage dit
beaucoup sur les actifs analysés et presque rien sur la part du parc qui
lui a échappé. Mme~EBODE présente au comité le nombre de failles
\emph{et} ces cinq lignes. Le premier chantier n'est pas la correction,
c'est l'inventaire.
\end{boxlecture}

% ============================================================================
\section{Ce que l'Identification demande à chaque Métier}
\label{sec:c10-metiers}
% ============================================================================

\begin{boxmetier}{Analyste en gestion des vulnérabilités}
\textbf{Identifier sans exploiter}
\begin{enumerate}
    \item Construire l'inventaire en croisant au moins trois sources.
    \item Préférer le balayage authentifié, et suivre le taux de réussite
          de l'authentification.
    \item Dédoublonner, puis vérifier un échantillon pour estimer les faux
          positifs.
    \item Exclure ou ménager les actifs fragiles, avec l'accord des métiers.
    \item Publier la couverture avec le nombre de failles.
\end{enumerate}
\end{boxmetier}

\begin{boxmetier}{RSSI}
Exiger de chaque prestataire qui installe un équipement sa déclaration à
l'inventaire, dans le contrat même. La passerelle de NYANGA et les 132
inconnus de SOLITRANS ont la même origine.
\end{boxmetier}

\begin{boxmetier}{Analyste en réponse à incident}
L'inventaire et l'historique des balayages servent aussi l'enquête~: ils
datent l'apparition d'une version vulnérable et montrent depuis quand une
faille était visible, ou invisible.
\end{boxmetier}

\separateur

\begin{boxresume}
\textbf{Ce qu'il faut retenir de ce chapitre~:}
\begin{itemize}
    \item L'identification commence par l'inventaire, construit en
          croisant plusieurs sources.
    \item Balayage non authentifié, authentifié, agent et analyse de
          configuration se complètent~; aucun n'exploite la faille.
    \item Les faux positifs se vérifient sur échantillon~; les faux négatifs
          se combattent par la couverture.
    \item Un balayage s'encadre~: autorisation écrite, périmètre, fenêtre,
          profil, contacts.
    \item Chez SOLITRANS, un actif sur onze était inconnu et 22~\% des
          serveurs analysés sans authentification.
\end{itemize}
\end{boxresume}

\begin{boxinfo}
\textbf{Points de contrôle.} Vous avez compris ce chapitre si vous savez~:
\begin{itemize}
    \item dire pourquoi un actif hors inventaire est un \NEK{}~;
    \item expliquer un correctif rétroporté et ses effets sur un rapport~;
    \item rédiger le cadre d'un balayage.
\end{itemize}
\end{boxinfo}

% ============================================================================
\section*{Exercices}
% ============================================================================

\begin{boxexo}[title={\ding{45}\quad Exercice~10.1. Quelle méthode~? \niveaubase}]
Choisissez une méthode d'identification pour~: (a) les portables des
commerciaux, rarement au bureau~; (b) les serveurs Linux du siège~; (c) le
portail des transitaires vu d'Internet~; (d) la configuration des
commutateurs~; (e) une application métier qui embarque des bibliothèques
open source.
\end{boxexo}

\begin{boxexo}[title={\ding{45}\quad Exercice~10.2. Faux positif ou pas \niveaumoyen}]
L'outil signale sur un serveur Debian une faille d'OpenSSH corrigée dans la
version 9.6, alors que le serveur annonce la version 9.2.
\begin{enumerate}[(a)]
    \item Pourquoi ce constat peut-il être un faux positif~?
    \item Comment le vérifier sans exploiter la faille~?
    \item Comment cotez-vous le constat avant et après vérification~?
\end{enumerate}
\end{boxexo}

\begin{boxexo}[title={\ding{45}\quad Exercice~10.3. La couverture \niveaumoyen}]
À partir de l'encadré « La couverture de SOLITRANS en mai 2025 »~:
\begin{enumerate}[(a)]
    \item Calculez le taux de couverture global (actifs analysés sur actifs
          existants).
    \item Si les 15~\% de faux positifs se confirment, combien des 2\,947
          constats seraient réels~?
    \item Pourquoi ce second chiffre est-il probablement trop bas~?
\end{enumerate}
\end{boxexo}

\begin{boxexo}[title={\ding{45}\quad Exercice~10.4. La passerelle invisible \niveauexpert}]
\begin{enumerate}[(a)]
    \item Décrivez trois contrôles qui auraient fait entrer la passerelle de
          NYANGA dans le périmètre du balayage.
    \item Une analyse non authentifiée de la passerelle, depuis Internet,
          aurait-elle détecté VULN-2024-4471~? À quelles conditions~?
    \item Rédigez le cadre (Tableau~10.4) d'un balayage mensuel de la
          zone réseau de la passerelle, installée et maintenue par un
          prestataire.
\end{enumerate}
\end{boxexo}

\bigskip

\begin{boxinfo}
\textbf{Pour aller plus loin.}
\begin{itemize}
    \item Le NIST SP~800-115~\cite{nist800115}, sur les techniques
          d'évaluation.
    \item Les CIS Benchmarks~\cite{cis2026benchmarks}, pour les références
          de configuration.
\end{itemize}
\end{boxinfo}

\begin{boxinfo}
\textbf{Transition.} SOLITRANS connaît désormais près de trois mille
failles, et sait qu'elle ne les connaît pas toutes. Lesquelles corriger
d'abord, dans quels délais, et comment suivre~? Chapitre~\ref{chap:gerer}.
\end{boxinfo}
